% Joint Bernoulli masses. Each row and column sums to 1/2.
\begin{tikzpicture}[x=1mm,y=1mm,
  every node/.style={font=\figurefont\footnotesize,text=FigureInk}]
  \foreach \offset/\title in {0/完全同向,38/独立,76/完全反向} {
    \begin{scope}[shift={(\offset,0)}]
      \draw[FigureStroke,line width=.6pt] (0,0) rectangle (28,24);
      \draw[FigureGrid,line width=.5pt] (14,0)--(14,24) (0,12)--(28,12);
      \node at (14,38) {\title};
      \node at (7,28) {\(V=0\)};
      \node at (21,28) {\(V=1\)};
      \node[rotate=90] at (-4,18) {\(U=0\)};
      \node[rotate=90] at (-4,6) {\(U=1\)};
    \end{scope}
  }
  \foreach \x/\y/\v in {7/18/0.5,21/18/0,7/6/0,21/6/0.5,
    45/18/0.25,59/18/0.25,45/6/0.25,59/6/0.25,
    83/18/0,97/18/0.5,83/6/0.5,97/6/0} {
    \node at (\x,\y) {\(\v\)};
  }
  \draw[FigureBlue,line width=1pt] (2,14) rectangle (12,22);
  \draw[FigureBlue,line width=1pt] (16,2) rectangle (26,10);
  \draw[FigureYellow,line width=1pt] (40,2) rectangle (64,22);
  \draw[FigureRed,dashed,line width=1pt] (92,14) rectangle (102,22);
  \draw[FigureRed,dashed,line width=1pt] (78,2) rectangle (88,10);
  \node[align=center] at (50,-8)
    {边缘相同：\(\mathbb P(U=1)=\mathbb P(V=1)=1/2\)};
\end{tikzpicture}
